\section{数值仿真实验与验证}
\label{sec:simulation}

基于第\ref{sec:dynamics}节的RK4高精度积分框架和第\ref{sec:INDI}节的INDI控制律方案，使用Python独立实现的六自由度仿真程序（\texttt{simulations/high-fidelity-analysis/}），对配平状态、线性化模态、积分精度和闭环响应进行定量分析与可视化。所有数值结果基于MODEL-DEFAULT参数，未经标定，结论限定为概念级。

\subsection{纵向配平}

在$V=24$\,m/s水平定直飞行条件下，Newton法联立求解纵向三自由度配平：

\begin{table}[H]
\centering
\caption{纵向配平状态 ($V=28\,\text{m/s}$，v0.2.0 VTOL 参数集，四方程)}
\label{tab:trim}
\small
\begin{tabular}{l c c}
\toprule
\textbf{参数} & \textbf{符号} & \textbf{值} \\
\midrule
迎角 & $\alpha$ & $4.122^\circ$ \\
尾摆偏置 & $\delta_{t0}$ & $-18.702^\circ$ \\
前摆偏置 & $\delta_{f0}$ & $-1.739^\circ$ \\
基准转速 & $\omega_0$ & $221.1\,\text{rad/s}$ \\
油门百分比 & — & $19.2\%$ \\
\bottomrule
\end{tabular}
\end{table}


该配平点（$\alpha=1.640^\circ$、$\delta_{t0}=-0.904^\circ$、油门$41.8\%$）与Web JS仿真（$\alpha_0=1.627^\circ$、$\delta_{t0}=-0.650^\circ$）存在小幅差异，来源为Python独立实现的简化气动模型。

\subsection{线性化与模态分析}

\begin{figure}[H]
\centering
\includegraphics[width=0.82\textwidth]{fig-sim/eigenvalues.pdf}
\caption{配平点特征值分布。蓝圈：纵向（短周期SP + 长周期Ph）；红方：横航向（荷兰滚DR + 滚转收敛RR + 螺旋Sp）。}
\label{fig:eigenvalues}
\end{figure}

\begin{table}[H]
\centering
\caption{配平点线化模态特征值 ($V=28\,\text{m/s}$，v0.2.0 参数集；线性化在三方程近似点 $\delta_f=0$ 处进行，$\delta_f=-1.739^\circ$ 的偏航耦合未计入)}
\label{tab:modes}
\small
\begin{tabularx}{\textwidth}{l c c c c}
\toprule
\textbf{模态} & \textbf{特征值} & \textbf{频率 (Hz)} & \textbf{阻尼比} & \textbf{时间常数/周期} \\
\midrule
  短周期 (SP) & $-3.365 \pm 5.509i$ & 1.027 & 0.521 & $\tau=0.297$\,s \\
  长周期 (Ph) & $+0.0177$ & 0.003 & $-1.000$ & 发散 $\tau=56.5$\,s（油门开环） \\
  航迹/高度 & $-0.170$ & 0.027 & 1.000 & $\tau=5.87$\,s \\
  滚转收敛 (RR) & $-47.65$ & 7.584 & 1.000 & $\tau=0.021$\,s \\
  荷兰滚 (DR) & $-1.425 \pm 6.861i$ & 1.115 & 0.203 & $T=0.916$\,s \\
  螺旋 (Sp) & $-0.156$ & 0.025 & 1.000 & $\tau=6.43$\,s \\
\bottomrule
\end{tabularx}
\end{table}


短周期模态（$f\approx2.6$\,Hz，$\zeta\approx0.84$）由$C_{m\alpha}$和$C_{mq}$支配，具备良好快速衰减；长周期（$f\approx0.05$\,Hz，$\zeta\approx0.35$）偏弱但SAS可改善。横航向：滚转收敛极快（$\tau\approx9.2$\,ms），荷兰滚（$f\approx2.0$\,Hz，$\zeta\approx0.48$）频率偏高——源于较小$I_z$和较大$C_{n\beta}$的组合，螺旋模态缓慢收敛（$\tau\approx8.5$\,s）。

\subsection{数值积分精度：RK4 vs 显式欧拉}

\begin{figure}[H]
\centering
\includegraphics[width=0.82\textwidth]{fig-sim/rk4_vs_euler.pdf}
\caption{RK4与显式欧拉的积分误差演化。RK4（$1.5\times10^{-12}$）远优于显式欧拉（$4.0\times10^{-5}$），精度提升$2.6\times10^7$倍。}
\label{fig:rk4_euler}
\end{figure}

恒定角速率$\omega=(0,0,1)^T$\,rad/s（纯偏航）、$dt=0.004$\,s、$T=60$\,s条件下，四元数范数偏差：$\lVert q_{\mathrm{Euler}}-q_{\mathrm{theo}}\rVert_2=4.0\times10^{-5}$，$\lVert q_{\mathrm{RK4}}-q_{\mathrm{theo}}\rVert_2=1.5\times10^{-12}$。图\ref{fig:rk4_euler}展示了两者误差随时间的演化趋势。RK4的误差已低于IEEE 754双精度浮点舍入累积效应，在60秒仿真中对姿态传播的影响可视为零，直接验证了第\ref{sec:dynamics}节的理论分析。

\subsection{SAS闭环响应}

\begin{figure}[H]
\centering
\includegraphics[width=0.88\textwidth]{fig-sim/sas_pitch_step.pdf}
\caption{SAS俯仰阶跃响应：上—俯仰角 $\theta$，中—俯仰角速率 $q$，下—尾摆指令 $\delta_t$。$t=2$\,s施加5\,deg/s的$q$脉冲。}
\label{fig:sas_response}
\end{figure}

$t=2$\,s时施加俯仰角速率脉冲（$\Delta q=5^\circ$/s），姿态偏差在约2秒内衰减至峰值的5\%以内（图\ref{fig:sas_response}），验证了角速率阻尼+姿态PI反馈的有效性。$\delta_t$峰值约$0.3^\circ$远低于$\pm25^\circ$限幅。

\begin{figure}[H]
\centering
\includegraphics[width=0.88\textwidth]{fig-sim/trim_stability.pdf}
\caption{配平稳定性：20秒SAS保持。上—前向速度 $u$，下—俯仰角 $\theta$。}
\label{fig:trim_stability}
\end{figure}

20秒配平保持（图\ref{fig:trim_stability}）空速漂移$<0.1\%$，表明SAS闭环在配平点附近维持稳定平衡。

\subsection{INDI控制器初步验证}

基于在线Jacobian $\mathbf{B}_{\mathrm{true}}$（式\ref{eq:B_true}）的INDI控制器在零扰动配平保持中表现与SAS相当（12秒速度漂移$<0.1\%$），验证了增量控制架构的配平稳定性。在角速率扰动响应中，INDI表现出比SAS更大幅度的执行器调节——这与INDI理论预测一致：增量控制通过$\mathbf{B}^{-1}$直接映射角加速度误差至执行器增量，对扰动的初始响应更激进。然而，在$\dot{\boldsymbol{\omega}}$估计仅依赖后向欧拉差分的纯传感器方案下，噪声放大效应导致INDI扰动抑制性能不如结构更保守的SAS。这一发现突显了第\ref{sec:INDI}节指出的角加速度滤波（§\ref{sec:accel_estimation}）对INDI实际性能的决定性影响。

\subsection{轴间耦合抑制：对角SAS与全矩阵INDI对比}
\label{sec:coupling_rejection}

第\ref{sec:allocation}节指出，前摆角通过反扭矩投影在俯仰通道产生交叉耦合$-\tau_f\sin\delta_f$（式\ref{eq:coupling_pitch}）。当前SAS（第\ref{sec:control}节）采用对角分配——每个通道独立映射到单一执行器，不含交叉项补偿，因此轴间耦合只能作为外部扰动由反馈被动抑制。INDI使用完整在线Jacobian $\mathbf{B}_{\mathrm{true}}$（式\ref{eq:B_true}）并全矩阵求逆，理论上可显式补偿该耦合。本节通过针对性数值实验定量对比两者的耦合抑制能力。

\textbf{实验设置}：飞行器处于24\,m/s巡航配平状态，$t=2$\,s时在前摆通道施加$\delta_f=10^\circ$阶跃（偏航指令），产生俯仰耦合扰动$-\tau_f\sin 10^\circ$。控制器需保持其余通道（俯仰/滚转/偏航角速率为零），测量俯仰角$\theta$偏离配平的峰值作为耦合抑制指标。陀螺噪声取标准差$0.002$\,rad/s的典型MEMS量级，每配置进行5次蒙特卡洛仿真取均值。INDI的角加速度估计采用后向欧拉差分。

\textbf{结果}如表\ref{tab:coupling_rejection}所示。

\begin{table}[H]
\centering
\caption{前摆$\delta_f=10^\circ$阶跃下的俯仰耦合抑制对比（5次蒙特卡洛均值）}
\label{tab:coupling_rejection}
\small
\begin{tabular}{l c c c}
\toprule
\textbf{控制架构} & \textbf{陀螺噪声} & \textbf{俯仰耦合峰值} & \textbf{稳态} $\boldsymbol{|\theta|}$ \\
\midrule
SAS（对角分配） & 0 & $1.63^\circ$ & $1.71^\circ$ \\
SAS（对角分配） & 0.002 rad/s & $1.63^\circ$ & $1.71^\circ$ \\
INDI（全矩阵$\mathbf{B}_{\mathrm{true}}$） & 0 & $\mathbf{0.03^\circ}$ & $0.03^\circ$ \\
INDI（全矩阵$\mathbf{B}_{\mathrm{true}}$） & 0.002 rad/s & $\mathbf{0.40^\circ}$ & $0.16^\circ$ \\
\bottomrule
\end{tabular}
\end{table}

结果揭示了两种架构的本质差异：

\begin{enumerate}[label=(\arabic*)]
    \item \textbf{SAS的耦合峰值（$1.63^\circ$）与噪声无关}——这是对角分配的\textbf{结构性}局限。由于分配矩阵不含交叉项，前摆产生的俯仰耦合$-\tau_f\sin\delta_f$完全未被前馈补偿，只能依赖俯仰角速率阻尼和姿态反馈被动衰减，其峰值由耦合量级与闭环带宽决定，与传感器噪声无关。这是几何解耦近似的固有代价（第\ref{sec:allocation}节）。
    \item \textbf{INDI在理想条件下几乎完全消除耦合（$0.03^\circ$）}——完整$\mathbf{B}_{\mathrm{true}}$的逆运算将前摆对俯仰的影响显式计入尾摆增量，实现了主动解耦。这验证了第\ref{sec:INDI}节的核心论点：INDI通过全矩阵Jacobian将构型的弱交叉耦合纳入控制律，而非视为扰动。
    \item \textbf{INDI的耦合抑制受角加速度估计质量制约}——噪声$0.002$\,rad/s下耦合峰值从$0.03^\circ$升至$0.40^\circ$。这是因为差分估计的噪声经$\mathbf{B}^{-1}$放大后污染增量指令。尽管如此，$0.40^\circ$仍优于SAS的$1.63^\circ$约4倍，表明即使在噪声条件下，全矩阵INDI的耦合抑制能力依然显著优于对角SAS。
\end{enumerate}

\textbf{结论}：对于本构型固有的反扭矩投影交叉耦合，全矩阵在线Jacobian的INDI相比对角SAS实现了数量级的耦合抑制改善（理想下约50倍，噪声下约4倍）。这是INDI相对SAS在控制质量上的\textbf{本质优势}，而非增量形式带来的边际改进。

\subsection{角加速度估计方法对INDI性能的影响}
\label{sec:estimator_study}

第\ref{sec:INDI}节§\ref{sec:accel_estimation}分析了六种角加速度获取方法。本节通过闭环蒙特卡洛仿真，定量评估其中四种可在线实现的方法对INDI耦合抑制性能的影响。实验场景同§\ref{sec:coupling_rejection}（前摆$\delta_f=10^\circ$阶跃），INDI内环带宽$K=3.0$，扫描陀螺噪声水平，每条件6次仿真取均值与标准差。

\textbf{估计器实现}（统一接口，接入INDI替换差分项）：
\begin{itemize}
    \item \textbf{后向欧拉差分}：基线，零延迟但噪声放大$\propto 1/\Delta t^2$；
    \item \textbf{Savitzky-Golay滤波}：窗口$M=3$、二阶多项式，群延迟约$4$\,ms；
    \item \textbf{固定滞后Kalman平滑}：滞后$N=2$步，延迟约$8$\,ms；
    \item \textbf{自适应互补滤波}：模型前向预测（含刚体转动、陀螺、Coriolis项）与数值差分融合，截止频率随角加速度残差自适应（$20$--$150$\,rad/s）。
\end{itemize}

\textbf{结果}如表\ref{tab:estimator_compare}所示（俯仰耦合峰值，均值$\pm$标准差）。

\begin{table}[H]
\centering
\caption{不同角加速度估计方法下INDI的俯仰耦合峰值（°，6次蒙特卡洛）}
\label{tab:estimator_compare}
\small
\begin{tabular}{l c c c c}
\toprule
\textbf{陀螺噪声} & \textbf{差分} & \textbf{S-G滤波} & \textbf{固定滞后平滑} & \textbf{自适应互补} \\
\midrule
0 & $\mathbf{0.13}$ & 2.81 & 3.18 & 1.77 \\
0.005 rad/s & 1.61$\pm$0.57 & 2.56$\pm$0.49 & 2.64$\pm$0.41 & 1.94$\pm$0.13 \\
0.01 rad/s & 2.92$\pm$1.30 & 2.70$\pm$0.49 & 3.06$\pm$0.83 & $\mathbf{2.02\pm0.33}$ \\
0.02 rad/s & 3.08$\pm$2.05 & 3.26$\pm$1.11 & 2.54$\pm$0.81 & $\mathbf{2.08\pm0.75}$ \\
\bottomrule
\end{tabular}
\end{table}

结果揭示了INDI角加速度估计的\textbf{延迟--噪声权衡}，以及一个反直觉但重要的结论：

\begin{enumerate}[label=(\arabic*)]
    \item \textbf{零噪声下差分最优（$0.13^\circ$）}：Savitzky-Golay与固定滞后平滑虽开环跟踪精度接近完美（幅值比$>0.995$），但其$4$--$8$\,ms群延迟在闭环中触发同步效应（第\ref{sec:INDI}节），$\nu-\dot{\boldsymbol{\omega}}$对消错位，耦合峰值反而升至$2.8$--$3.2^\circ$——\textbf{比不滤波更差}。这表明INDI估计器的评价标准不是开环精度，而是闭环同步性。
    \item \textbf{中高噪声下自适应互补最优且最稳健}：噪声$\ge 0.01$\,rad/s时，自适应互补滤波的耦合峰值（$\sim 2.0^\circ$）与标准差（$\pm 0.33$）均为最小。其模型预测通道提供低噪声低频分量、差分通道保留高频低延迟，绕开了纯滤波器的延迟--噪声权衡。这与Jiang等\cite{jiang2025adaptive}的结论一致，并首次在本构型上定量验证。
    \item \textbf{自适应互补的优势依赖模型通道准确性}：进一步的完整闭环实验发现，当模型预测通道所用的控制输入先验与实际电机响应存在偏差（电机一阶滞后导致增量指令与实际推力不一致）时，自适应互补的估计方差反而增大（$\pm 0.95$）。因此其工程应用前提是模型通道与实际执行器动态的一致性标定。
\end{enumerate}

\textbf{工程建议}：概念仿真阶段（理想$\dot{\boldsymbol{\omega}}$）无需滤波；HIL与实机阶段（噪声$\ge 0.005$\,rad/s）推荐自适应互补滤波，但须先标定模型预测通道与执行器实际响应的一致性；不推荐固定滞后平滑，其延迟代价在INDI闭环中超过噪声收益。

\subsection{模型参数敏感性}

上述数值结果对气动导数不确定性敏感。以短周期模态频率为例，若$C_{m\alpha}$偏离MODEL-DEFAULT值$\pm30\%$，短周期频率从2.6\,Hz移至2.2--3.0\,Hz，阻尼比从0.84变为0.72--0.93。该量级变化在常规SAS增益补偿能力之内，但突显了台架标定和CFD验证对飞行品质定量评估的必要性。
